% P1 the goal -> P2 prior work trains a VAE, and validates on synthetic data ->
% P3 the failure we measured (z_dyn absorbs everything; the standard penalty cannot
% fix it) -> P4 the structural fix + rate -> P5 what it buys -> contributions -> scope.
% Numbers: DEVLOG v6.2. Claims dropped in this revision are noted at the end of the file.

A video representation should keep its kinds of information in different places: what a
scene \emph{contains}, and how it \emph{moves}. Video VAEs do the opposite. Trained to
reconstruct, they pack texture, identity and motion into a single unstructured latent
grid offering no axis along which a downstream user can read out identity, hold it
fixed, or edit it \citep{videomae,vjepa}. That structure can be recovered, but
expensively: sequential VAEs that split content from motion do so by \emph{training a
VAE from scratch} \citep{s3vae,cdsvae,vidtwin}, recent predictive video VAEs enrich the
latent but leave it entangled \citep{pvvae}, and diffusion-based factorizations train
the generator end-to-end \citep{divid}. We ask a cheaper question: given an
\emph{already-trained, frozen} video VAE, can we re-encode its latent into a static code
and a dynamics code, at high compression, without retraining it?

The obvious construction --- two encoder trunks, two codes, one decoder reading both ---
fails in a way that is easy to miss and, once measured, hard to unsee. \textbf{The
dynamics code absorbs the entire scene.} It is paid once per frame and is therefore the
larger code, so a decoder free to read both simply uses it for everything. In our
setting, deleting the static code raises reconstruction error by only $11\%$ while
deleting the dynamics code raises it by $714\%$; the static code decodes to $24.90$\,dB
on its own, and giving it eight times the rate moves that to $24.89$\,dB while leaving
three quarters of its dimensions unused. The split exists on paper and not in the model.

The standard remedy --- penalising dependence between the codes, as a
cross-correlation or an orthogonality term --- cannot repair this, for a reason that
holds regardless of implementation: \textbf{a static code carrying pure noise satisfies
it perfectly}, because noise is uncorrelated with everything. The penalty discourages
redundancy; it is silent on usefulness. We confirm this directly. In the failing model
above the penalty is essentially fully satisfied ($3.7\times 10^{-3}$) while the static
code holds almost nothing, and across models whose measured code overlap spans
$0.32$--$0.51$ (kernel CKA) the penalty reads $0.0028$--$0.0053$ --- it does not rank
them. Removing it improves reconstruction by $21\%$ at no cost to attribute readout.

So we stop asking the loss to enforce the split and let the decoder enforce it.
\method{} reconstructs each frame as a static image plus a \emph{zero-mean} per-frame
residual, where the residual branch never sees the static code and its temporal mean is
subtracted after the nonlinearity. The time-constant part of the output is then produced
by the static code \emph{identically}, not approximately, and the dynamics code can only
describe deviation from it. Deleting the static code now costs $484\%$ rather than
$11\%$. We pair this with a rate correction the diagnosis implies: $91.5\%$ of a chunk's
latent energy is time-constant, yet the conventional allocation spends $4\%$ of its
floats there, paying for static content once per frame instead of once per chunk.

Rebalancing the budget makes the model \emph{smaller}. At $576$ static floats and $64$
per frame --- $1{,}152$ floats per chunk, a $24\times$ compression and half the rate of
our own prior configuration --- \method{} matches VideoFlexTok and VideoMAE on
reconstruction under a matched protocol, while carrying a factorization neither
expresses. On $163{,}717$ Something-Something-v2 clips the dynamics code reads actions
at $10.3\%$ against the static code's $6.1\%$ (chance $1.9\%$), which is what a working
static/dynamic split predicts on a motion benchmark.

Our contributions are:
\begin{itemize}
  \item \textbf{A diagnosis} (\cref{sec:absorb}): in a two-code video autoencoder the
  per-frame code absorbs the scene, and additional capacity for the static code does not
  help --- eight times the rate changes its solo reconstruction by $0.01$\,dB.
  \item \textbf{A negative result about the standard fix} (\cref{sec:indep}):
  independence penalties are satisfied by a noise static code and, measured against
  kernel CKA, fail to order models by their actual code overlap. We drop the term.
  \item \textbf{A structural fix} (\cref{sec:basedelta}): a base-plus-zero-mean-residual
  decoder makes the time-constant output provably a function of the static code alone,
  raising its ablation cost from $11\%$ to $484\%$.
  \item \textbf{A rate correction} (\cref{sec:rate}): matching the bit allocation to
  where the information is yields $24\times$ compression at parity with strong video
  encoders under a matched protocol.
  \item \textbf{The spatial-dynamics finding} (\cref{sec:spatial-dyn}): a global
  per-frame dynamics code is rank-limited and fails on real video; a per-frame spatial
  grid repairs it ($+4.4$\,dB), attributable to structure rather than rate.
\end{itemize}

We are explicit about scope. \method{} trades absolute fidelity for structure: at
matched settings it reconstructs $2.3$\,dB below an unconstrained model of the same
design, and it is not a codec. Its semantic accuracy is bounded by the frozen latent it
distils, not by the encoder --- a linear probe on the \emph{full} latent reaches
$13.7\%$ on Something-Something-v2 against our dynamics code's $10.3\%$ at $48\times$
fewer floats. And we report one negative that bears on how such claims should be
evaluated at all: an identity/motion swap test, standard in this literature, is passed
just as well by an encoder with a \emph{shared} trunk that cannot factorize
(\cref{sec:q2}) --- it measures the decoder's wiring, not the encoder's split, and needs
that control to be evidence.

% DROPPED IN THIS REVISION, and why:
%  - "cleanly disentangled (diagonal-dominant cross-probe)" and the identity/motion swap
%    as evidence: the swap is passed by a shared-trunk control (+0.054 vs +0.053+-0.013),
%    so it does not separate a factorizing encoder from an entangled one.
%  - "z_static beats raw-latent, PCA and random-init at 288x smaller": contradicted by
%    our own Table 1, where the raw latent reads 0.785 colour mAP against z_static's
%    0.767. The abstract's "within noise of the full latent" was the correct phrasing.
%  - the camera-pose path: unchanged in the code but not part of the committed
%    configuration; it belongs in the moving-camera section, not the contributions.
